% Transcription of folder 149, batch 06 (archivists' pages 101-120).
% Pass: Opus 5.5 (claude-opus-5-5), 2026-10-03 — first pass, unchecked against the pages by a human.
%
% Pages skipped: none. Grothendieck's own pagination (51-60) is noted per page.
\documentclass[11pt,a4paper]{article}
\input{../preamble/grothendieck}

\folder{149}
\batch{6}
\pages{101}{120}
\foldertitle{[Autour de La "Longue Marche" à travers la théorie de Galois, pages 1 à 67] : notes manuscrites (s.d.).}
\dating{s.d. — le groupe « Autour de La "Longue Marche" à travers la théorie de Galois » (141 à 149) est daté [à partir de 1978]-1983}
\watermark{Édition de démonstration}

\begin{document}

\note{notations du lot : $\mathcal{E}$ rend la lettre ronde par laquelle l'auteur désigne le groupe fondamental profini (« $\mathcal{E}_X = \pi_1(X)$ », p.~104), l'exposant marquant le revêtement universel ou le point géométrique choisi ; $\pi$ garde le $\pi$ de l'auteur}

\page{101}
\note{la phrase qui ouvre la page commence au feuillet précédent}
(Ce qui signifie aussi que $H$ (\uncertain{seul} \add{(\ill{} \ill{})} \uncertain{dans} \ill{}
de 31 --- p.\ ex.\ défini par un pt de $U_{\bar{\xi}}$
\ill{} $/K$ --- \uncertain{induit} un homom.\ $T \to \pi$)

c) L'action de $T$ sur $\pi_{\mathrm{ab}} = H_1(U_{\bar{\xi}})$ est
triviale.

d) Item pour $H_1(U_{\bar{\xi}}, T_\ell)$, où $\ell$ est premier
à la car.\ résiduelle. (\textbf{NB} \struck{\ill{}} \add{$\ell = p > 0$}
\uncertain{on} doit pouvoir l'exprimer par un critère
sur le groupe $p$-divisible…)

e) (\uncertain{Etant} \uncertain{donnée} une section de (30),
définissant une extension de $\mathcal{E}_k$ par $\pi$,
et supposant maintenant $k$ de type
fini sur $\mathbb{Q}$) La dite extension est
admissible (\uncertain{quitte} à passer à une
\add{\uncertain{ext.}\ de $\mathcal{E}_k$,} \ill{} \ill{} \uncertain{sous-groupe} \ill{} \uncertain{ouvert}), i.e.\
\ill{}, \uncertain{toujours} \ill{} peut être défini par une courbe alg.\
\uncertain{sur} \uncertain{une} \add{$k$, \ill{} \uncertain{sur}} extension finie de $k$.

f) Cette \uncertain{dernière} \uncertain{dans} \ill{} d'une section de (30)
et d'une de (31)., \uncertain{permettant} de
faire \uncertain{opérer} \add{\uncertain{par} \ill{}} $\mathcal{E}_k$ sur $\pi$) : Cette
action \uncertain{doit} \uncertain{être} \uncertain{(remarque} (où \uncertain{définition} \uncertain{commence}
\uncertain{seulement} !)
\note{les renvois (30) et (31) visent des formules antérieures au lot}

\page{102}\note{p. 51 de l'auteur}
On a évidemment a) $\Rightarrow$ b) $\Rightarrow$ c) $\Rightarrow$ d),
et d) $\Rightarrow$ a) devrait être plus ou moins
connu, au moins de \uncertain{prouver} à
l'aide de choses bien connues. Ces
conditions impliquent e), mais je n'ai
pas d'idée pour le moment comment
voir que e) implique bonne réduction ---
il faudrait une compréhension (prof. ou
\uncertain{arithmétique}) plus approfondie que celle
que j'ai pour le moment.
\marginal{Ça pourrait se \uncertain{faire} \uncertain{en} \uncertain{réduisant} au cas où $\pi \simeq$ \ill{} \ill{} $+\, y$ \ill{} \ill{} \ill{} \ill{} $+$ \ill{} \ill{} \ill{} $S$ \ill{} \ill{} \ill{} \ill{} \ill{} \uncertain{devrait} \ill{} \ill{} \ill{}…}
\note{marginale écrite en oblique dans la marge gauche, très serrée ; seuls quelques mots se lisent}

J'ai donc là un principe général
(que je \uncertain{pourrais} à l'occasion formuler
avec la généralité qui lui \uncertain{revient}) pour
\add{(\uncertain{± naturelles})} \ill{}
construire des actions extérieures
de groupes $\mathcal{E}_k$ ($k$ ext.\ de type fini
de $\mathbb{Q}$) sur des $\pi$ à \uncertain{lacets}, qui
(sans doute…)
ne soient \uncertain{pas} géométriques, par \struck{des} la
considération de courbes de type gr
« se réduisant mal ». Mais je n'arrive
\note{la phrase se poursuit à la page suivante}

\page{103}
pas par cette voie à faire des actions
effectives, ou \uncertain{associées} à \struck{\ill{}} des actions
ext.\ déjà géométriques, qui ne soient
elles-mêmes géométriques --- i.e.\ d'une
des deux types $1^{\circ}$, $2^{\circ}$ envisagés au
début de ce \uncertain{n°}.
\note{le renvoi « $1^{\circ}$, $2^{\circ}$ » vise un passage antérieur au lot ; le reste du feuillet est blanc}

\page{104}\note{p. 52 de l'auteur}
\section{4. Sections d'extensions, et anneaux de valuation généraux.}
\note{titre souligné par l'auteur ; « d'extensions, » est ajouté par lui au-dessus de la ligne}

D'abord une révision des notations. Si
$X$ est un schéma connexe, je \uncertain{dénote} \struck{$\mathcal{E}_X$}
\[
\text{(1)} \qquad \mathcal{E}_X = \pi_1(X)
\]
son groupe fond.\ profini en tant que
groupe extérieur, et si $\tilde{X}$ est un
rev.\ universel profini de $\mathcal{E}_X$, par
\[
\text{(2)} \qquad \mathcal{E}^{(\tilde{X})}_{X} = \pi_1(X ; \tilde{X}) = \operatorname{Aut}_X(\tilde{X})
\]
son groupe fondamental précisé --- qui est un
groupe profini. Si $X = \operatorname{Spec}(A)$, où $A$
est un anneau (le plus souvent
un corps) je note $\mathcal{E}_A$, et $\mathcal{E}^{\tilde{A}}_{A}$ (si $A$
\add{$= \mathcal{E}^{\tilde{X}}_{X}$}
est hensélien…) \uncertain{(où} \uncertain{ces} \ill{} \uncertain{telle} que $\operatorname{Spec}(\tilde{A}) = \tilde{X}$
soit un rev.\ universel de $X$ (ce qui
le détermine à \uncertain{isom.}\ unique
près). Bien entendu, si $\xi$ est un
« point géométrique » de $X$, on note
\[
\text{(3)} \qquad \mathcal{E}^{\xi}_{X} = \pi_1(X, \xi) = \mathcal{E}^{\tilde{X}(\xi)}_{X} ,
\]
où $\tilde{X}(\xi)$ est le \uncertain{rev.}
universel de $X$ pointé en $\xi$. Le choix
\note{« (si $A$ est hensélien…) » : lecture très incertaine du mot rattaché à $A$}

\page{105}
de $\xi$ correspond au choix d'un pt $x$ de $X$
et d'une ext.\ sép.\ close $\overline{\Omega}$ de $k(x)$
(par ex.\ une clôture alg.\ de $k(x)$) et on
note \uncertain{aussi} \quad $\mathcal{E}^{\Omega}_{X}$ --- \uncertain{Lieu} de
\uncertain{notation} \uncertain{au} \ill{} \ill{} \uncertain{qu'}$\Omega$
\uncertain{extension} de $k(x)$, dans \uncertain{un cas} \uncertain{abstraction}
\add{\uncertain{\ill{} \ill{}}}
de $k(x)$ \uncertain{algébrique}). \uncertain{On} \uncertain{a}
\[
\text{(4)} \qquad \mathcal{E}^{\Omega}_{x} \simeq \mathcal{E}^{\overline{k(x)}}_{X} ,
\]
où $\overline{k(x)}$ est la clôture sép.\ de $k(x)$ dans
$\Omega$. Bien sûr, si $X = \operatorname{Spec}(A)$, on
notera aussi \quad $\mathcal{E}^{\Omega}_{A}$ --- notation \uncertain{standard}
utilisée pour \quad $\mathcal{E}^{\overline{K}}_{K}$, $K$ un corps,
$\overline{K}$ une clôture alg.\ ou sép.\ de $K$.

Si $X$ est un $Y$-schéma, \struck{on \uncertain{notera}} \add{$Y$ connexe,}
\[
\text{(5)} \qquad \mathcal{E}(f) : \mathcal{E}_X \longrightarrow \mathcal{E}_Y \qquad \text{où } f : X \to Y
\]
($\mathcal{E}_X$ est un foncteur en $X$
(Schémas connexes) $\longrightarrow$ groupes ext.
\note{dans (4), l'indice du premier membre est écrit $x$ (ou $\alpha$) et non $X$}

\page{106}\note{p. 53 de l'auteur}
qui se précise par l'homom.\ effectif de groupes
profinis
\[
\text{(6)} \qquad \mathcal{E}^{\tilde{X}}_{X} \longrightarrow \mathcal{E}^{\tilde{Y}}_{Y}
\]
où $\tilde{Y}$ est le rev.\ universel de $Y$ défini
par $\tilde{X} \to Y$ ($\tilde{X}$ jouant le rôle de
\uncertain{point} --- \struck{\ill{}} \uncertain{devrait} pouvoir écrire en
fait $\mathcal{E}^{\tilde{X}}_{Y}$, plus généralement $\mathcal{E}^{Z}_{Y}$ chaque fois
qu'on a un $Y$-schéma $Z$ 1-connexe,
jouant le rôle de « fonction fibre » pour
le topos $B_{\pi(X)}$ des rev.\ étales de $Y$…)
\struck{\ill{} $\mathcal{E}^{\alpha}_{X}$}
On \uncertain{peut} \uncertain{dire} que
\[
\text{(7)} \qquad (X, \tilde{X}) \longmapsto \mathcal{E}^{\tilde{X}}_{X}
\]
est un foncteur, de la catégorie des
schémas connexes $X$ \emph{munis} d'un
rev.\ universel (ou seulement d'un
$Z$ 1-connexe s'envoyant sur $X$) vers
celle des groupes profinis. Ceci s'applique
\add{\uncertain{\ill{}}}
en particulier en regardant la cat.\
des $(X, \xi)$ munis d'un pt géom.\ ---
\note{la phrase se poursuit à la page suivante}

\page{107}
on a donc
\[
\text{(8)} \qquad \mathcal{E}^{\xi}_{X} \longrightarrow \mathcal{E}^{\eta}_{\text{\struck{$Y$}}\,Y}
\]
\note{l'indice du second membre est un $Y$ biffé puis récrit ; l'exposant est lu $\eta$ sans certitude}
si $\ldots$ un homom.\ de schémas « géométriquement
pointés » $(X, \xi) \to (Y, \eta)$. On
\uncertain{note} \uncertain{que} \uncertain{la} \uncertain{ponctuation} géométrique
de $X$ en un $x \in X$ --- i.e.\ une
ext.\ \uncertain{sép.}\ close $\Omega$ de $k(x)$ ---
définit \add{\uncertain{canoniquement}} une ponctuation géométrique
de $Y$ en $y = f(x)$, \struck{\ill{}} par composition
--- $\Omega$ « est » une extension de $k(y)$ puisque
$k(x) \hookleftarrow k(y)$ --- et l'homom.\ (8)
\note{sous $k(x)$, un petit $x$ récrit}
s'identifie aussi à
\[
\text{(9)} \qquad \mathcal{E}^{\overline{k(x)}}_{X} \longrightarrow \mathcal{E}^{\overline{k(y)}}_{Y} ,
\]
où $\overline{k(x)}$, $\overline{k(y)}$ sont les clôtures séparables
\emph{dans $\Omega$}.

On posera
\[
\text{(10)} \qquad \mathcal{E}^{\tilde{X}}_{X/Y} = \operatorname{Ker}\bigl(\mathcal{E}^{\tilde{X}}_{X} \longrightarrow \mathcal{E}^{\tilde{X}}_{Y} = \mathcal{E}^{\tilde{Y}}_{Y}\bigr)
\]

\page{108}\note{p. 54 de l'auteur}
(C'est un foncteur par rapport au triple
$\tilde{X} \to X \to Y$ avec $X, Y$ 0-connexes,
$\tilde{X}$ un rev.\ universel ; plus généralement,
\struck{$Z$} \add{$T$} si $T \to X$ avec $T$ 1-connexe,
\uncertain{on} pose
\[
\text{(11)} \qquad \text{\struck{$\mathcal{E}$}}\,\mathcal{E}^{T}_{X/Y} = \operatorname{Ker}\bigl(\mathcal{E}^{T}_{X} \longrightarrow \mathcal{E}^{T}_{Y}\bigr)
\]
\note{dans (11) et dans la ligne qui précède, $T$ est écrit sur un $Z$ biffé}
fonctoriel en
\add{\uncertain{(lisible)}}
$(T \to X \to Y)$, \uncertain{en particulier} : $\mathcal{E}^{\xi}_{X/Y}$,
\uncertain{si} $\xi$ un pt géom.\ de $X$, $\mathcal{E}^{\Omega}_{X/Y}$,
$\mathcal{E}^{\tilde{X}}_{X/A}$ (si $Y = \operatorname{Spec} A$), $\mathcal{E}^{\tilde{B}}_{B/A}$…
\marginal{\textbf{NB} $\mathcal{E}^{T}_{X/Y}$, (11) ! \ill{} $X/Y$ \ill{} \ill{} \ill{} \ill{} \ill{} \ill{} \ill{} \ill{} \ill{} \ill{} \ill{} $\mathcal{E}_{X/Y}$}
\note{marginale oblique au milieu de la marge gauche, très serrée}

Un cas intéressant est celui où on
dispose d'une « suite exacte d'homotopie universelle »
\struck{\ill{}} (en dim 1)
pour $X \to Y$, alors la donnée
(par exemple $X = Y$ donné) de $T \to X$,
\marginal{Car \uncertain{$\mathcal{E}^{T}_{\ldots} = \mathcal{E}^{T}_{Y}$} \ill{} \uncertain{épi}…}
\uncertain{revient} $T$ 1-connexe) peut s'interpréter
par la donnée du composé
\note{la phrase se poursuit à la page suivante}

\page{109}
\[ T \longrightarrow Y \]
\emph{et} d'un relèvement de $T \to X$, ce qui \uncertain{revient au même}, d'une section
de $X_T = X \times_Y T$ sur $T$. Ceci posé,
on \uncertain{a} un \uncertain{isom.}
(comme il \uncertain{résulte} de l'hyp.\ de « suite exacte d'homotopie » \uncertain{faite})

\begin{tikzcd}
X \arrow[d] & X_T \arrow[l] \arrow[d] & \\
Y & T \arrow[l] & T \arrow[ul] \arrow[l, no head, "="]
\end{tikzcd}

\[
\text{(12)} \qquad \mathcal{E}^{T}_{X/Y} \simeq \pi_1(X_T ; T) \simeq \mathcal{E}^{T}_{X_T}
\]
et \uncertain{une suite exacte}
\[
\text{(13)} \qquad 1 \longrightarrow \mathcal{E}^{T}_{X/Y} \longrightarrow \mathcal{E}^{T}_{X} \longrightarrow \mathcal{E}^{T}_{Y} \longrightarrow 1
\]
(Cette hypothèse \uncertain{se trouve satisfaite} si $Y = \operatorname{Spec} K$, $K$ un corps,
si $X$ est géom. connexe sur $K$.)
Plus généralement, si \uncertain{on a une} suite exacte d'homotopie universelle,
mais pas \uncertain{univ.} \ill{} $\mathcal{E}^{T}_{X} \to \mathcal{E}^{T}_{Y}$ \uncertain{surj.}),

\page{110}\note{p. 55 de l'auteur}
\marginal{dans l'hyp. d'« homo\ill{} » \ill{} « suite exacte » \ill{} \uncertain{fibres} \ill{} topie », mais \ill{} pas \uncertain{univ.}}
on a une factorisation \uncertain{canon.} de $X \to Y$ en
\[
\text{(14)} \qquad X \longrightarrow Y' \longrightarrow Y
\]
avec $Y'$ étale fini \uncertain{ou} pro-étale fini sur $Y$
et $\mathcal{E}_X \to \mathcal{E}_{Y'}$ étant maintenant épi,
\struck{donc} cette suite exacte univ. d'homotopie \uncertain{bien vérifiée}. On a \ill{} dans ce cas
\[
\text{(15)} \qquad \mathcal{E}^{T}_{X/Y} \xrightarrow{\ \sim\ } \mathcal{E}^{T}_{X/Y'} ,
\]
ce qui peut donc se calculer comme
\[
\text{(16)} \qquad \mathcal{E}^{T}_{X \times_{Y'} T} \simeq \mathcal{E}^{T}_{X/Y}
\]
\note{la formule porte bien $X/Y$ à droite, et non $X/Y'$}
qui peut \uncertain{s'interpréter} \ill{} $\mathcal{E}^{T}_{X_T}$, mais \uncertain{en}
\struck{\ill{}} interprétation \ill{} \ill{} $X_T$ désigne non plus $X \times_Y T$ (qui
\ill{} à $Y' \to Y$ \uncertain{pas} \ill{}) mais $X \times_Y T$.
\note{dans ce dernier $X \times_Y T$, le premier $X$ est repassé à l'encre forte ; l'indice attendu serait $Y'$}
Bien sûr, à isom. \uncertain{près}, \ill{} pour \ill{} \struck{\ill{}} \uncertain{remplace}, \ill{}
le contenu de $\mathcal{E}^{T}_{X_T}$,
$T$ par $\tilde{Y}$ pour un \uncertain{point} \ill{} de $X_T$, et par
$\tilde{X}$ pour ce qui est de l'analogue des $X_{\tilde{Y}}$,

\page{111}
où $\tilde{X}, \tilde{Y}$ sont les rev. universels de $X, Y$ déduits de
\struck{\ill{}} \ill{} $\overline{T} \to Y$, \ill{} qu'on \uncertain{obtient})
\[
\text{(17)} \qquad \mathcal{E}^{T}_{X/Y} \simeq \mathcal{E}^{\tilde{X}}_{X/Y} \simeq \operatorname{Ker}\bigl(\mathcal{E}^{\tilde{X}}_{X} \to \mathcal{E}^{\tilde{Y}}_{Y}\bigr)
\]
(sous l'hyp. de suite exacte) et
\[
\text{(18)} \qquad \mathcal{E}^{T}_{X/Y} \simeq \mathcal{E}^{Z}_{X_{\tilde{Y}}} \simeq \mathcal{E}^{\tilde{X}}_{X_{\tilde{Y}}}
\]
\note{la formule (18) est écrite au-dessus d'une première version, entièrement biffée et illisible, qui portait encore $\mathcal{E}^{\tilde X}_{X_{\tilde Y}}$ ; la fin $\mathcal{E}^{\tilde X}_{X_{\tilde Y}}$ porte elle-même une surcharge}
dans le cas général.

\begin{tikzcd}
X \arrow[d] & X_{\tilde{Y}} \arrow[l] \arrow[d] & \tilde{X} \arrow[l] & Z \arrow[l] \\
Y' \arrow[d] & \tilde{Y} \arrow[l] & & \\
Y & & &
\end{tikzcd}

Une \ill{} des \ill{} de $\mathcal{E}^{T}_{X/Y}$ s'obtient \uncertain{en supposant} \struck{\ill{}}
$T$ \uncertain{réduit} à un pt géométrique, qui définit $\xi_X$, un pt géom. $\xi_Y$ de $Y$,
une fibre géom. $X_{\xi_Y}$ de $X$ sur $Y$ \ill{} $\xi_Y$, $\xi_X$ s'interprète comme un
pt géom. de cette fibre, et on a
\[
\text{(19)} \qquad \mathcal{E}^{\xi_X}_{X/Y} \simeq \mathcal{E}^{\xi_X}_{X_{(\xi_Y)}} \simeq \pi_1\bigl(X_{(\xi_Y)}, \xi_X\bigr)
\]
--- i.e. $\mathcal{E}_{X/Y}$ s'interprète comme groupe fondamental
des « fibres géom. » de $X/Y$.

\page{112}\note{p. 56 de l'auteur}
La \ill{} \ill{} \struck{définition} \add{\uncertain{les conventions}} (11)
et décidons que, si $\xi_X$ provient d'un pt géom. $\xi_T$ de $T$, alors on aura
des isos canon.
\[
\begin{aligned}
\text{(20)} \qquad \mathcal{E}^{\xi_T}_{X/Y} & \simeq \mathcal{E}^{T}_{X/Y} \\
\wr\ \ & \\
\pi_1\bigl(X_{\xi_Y}, \xi_T\bigr) &
\end{aligned}
\]
\struck{\ill{}} --- les $\pi_1(X_{\xi_Y}, \xi_T)$ des fibres
géom., relatifs aux divers pts géom.
de $\mathcal{E}^{T}_{X/Y}$, sont \uncertain{canoniquement} isom.
entre eux.
\note{« pts géom. de $\mathcal{E}^{T}_{X/Y}$ » : c'est bien ce que porte la page ; on attendrait les points géométriques $\xi_T$ de $T$}

(Il faut bien faire attention (\uncertain{$Z$ non précisé})
que $\mathcal{E}_{X/Y}$ \uncertain{n'est} une classe d'isom. de groupes profinis \uncertain{sans} plus,
\ill{} un groupe extérieur. L'arbitraire de la donnée d'un \uncertain{isom.} entre
$\mathcal{E}^{T}_{X/Y}$ et $\mathcal{E}^{T'}_{X/Y}$ provient des choix

\page{113}
d'un \ill{} de $\mathcal{E}^{T}_{X}$ \struck{$/Y$}, \uncertain{opérant}
sur le \ill{}-groupe $\mathcal{E}^{T}_{X/Y}$ par \uncertain{auto.} intérieurs.
\marginal{\ill{} relation \ill{} $\pi^{T}_{X/Y}$ ; que \ill{} \ill{} \ill{} \ill{} \ill{} \uncertain{temporairement}.}

Les relations ($\mathcal{E}_X$, $\mathcal{E}_{X/Y}$ \uncertain{nouvelles}) \uncertain{relativement} \ill{} pt de vue qui
regarde $\mathcal{E}^{T}_{X}$ comme une \emph{extension}
de $\mathcal{E}^{T}_{Y}$ (ou d'un sous-groupe \uncertain{convenable}) par $\mathcal{E}^{T}_{X/Y}$,
qui \ill{}, pour les besoins de la cause, \ill{} $\pi$ --- il joue le
rôle d'un groupe fondamental (géométrique)), tandis que $\mathcal{E}^{T}_{Y}$ joue
le rôle d'un groupe de Galois --- d'un groupe profini \ill{}.

Le cas le plus important pour nous
sera celui où $Y = \operatorname{Spec} K$, $K$ un corps.
\struck{Aussi formellement} Dans ce cas, $\mathcal{E}_{Y} \simeq \mathcal{E}^{\overline{Y}}_{Y}$

\page{114}\note{p. 57 de l'auteur}
s'identifie bel et bien au groupe de Galois $\operatorname{Gal}(\overline{K}/K)$ de \struck{\ill{}} $\overline{K}/K$,
$\overline{K}$ étant la clôture sép. de $K$ telle que $\operatorname{Spec} \overline{K} \simeq \overline{Y}$.
Souvent, on notera $\Gamma_K$, $\Gamma^{\overline{K}}_{K}$, $\Gamma$, (au lieu de $\mathcal{E}_Y$) --- \ill{}

si $K$ est alg. sur le corps premier, \uncertain{et qu'il} \ill{} donc plus de « partie
géométrique » à distinguer d'une « partie arithmétique »…

Soit $K$ \struck{\ill{}} \add{\ill{}} corps (qui pourrait être alg. clos), $L$ une \uncertain{ext.} de
type fini de $K$, $X$ un « modèle » propre régulier de $L$. Alors
$\mathcal{E}^{\overline{L}}_{X}$ s'identifie à un quotient de $\mathcal{E}^{\overline{L}}_{L}$,
qui \emph{ne dépend pas du modèle $X$ choisi}, comme il est bien connu
\note{la page porte « des modèles / $X$ choisis » ; le passage souligné l'est par l'auteur}

\page{115}
\marginal{$\mathcal{E}^{\mathrm{g}\overline{L}}_{L}$ \ill{} \struck{\ill{}} \ill{} $\mathcal{E}^{\overline{L}}_{L}$}
c'est la partie « universellement générique » de
$\mathcal{E}^{\overline{L}}_{L}$ \struck{relative à \ill{} qui \ill{} les \ill{}} (qui \ill{} \ill{})
classifie les \uncertain{sel.} (finis) \ill{} sur $L$
qui sont « \uncertain{non ramifiés} » sur tout modèle
régulier (propre ou non) de $L/K$.
\note{la petite lettre placée devant $\overline{L}$ dans l'exposant, ici et en (21), (23), est lue $\mathrm{g}$ sans certitude}
Si $U$ est un modèle quelconque, il se plonge dans les $X$, et on a des homomorphismes surjectifs
\[
\text{(21)} \qquad \mathcal{E}^{\overline{L}}_{L} \longrightarrow \mathcal{E}^{\overline{L}}_{U} \longrightarrow \mathcal{E}^{\overline{L}}_{X} = \mathcal{E}^{\mathrm{g}\overline{L}}_{L}
\]
On aura \struck{\ill{}}
\[
\text{(22)} \qquad U = X \setminus Z
\]
$Z$ partie fermée de $X$ --- soient $D_i$ ses composantes de codim 1, $x_i$ les
pts génériques des $D_i$. On a des
groupes de décomposition et d'inertie dans $\mathcal{E}^{\overline{L}}_{U}$
\[
N_{x'_i} \quad \text{et} \quad I_{x'_i}
\]
(\struck{d} relatifs au choix d'un $x'_i$ au dessus de $x_i$ dans le \uncertain{normalisé} de
$X$ dans $\overline{L}$), et on trouve

\page{116}\note{p. 58 de l'auteur}
\[
\text{(23)} \qquad \mathcal{E}^{\overline{L}}_{X} = \mathcal{E}^{\mathrm{g}\overline{L}}_{L} \simeq \mathcal{E}^{\overline{L}}_{U} \big/ \text{ss-groupe fermé invariant engendré}
\]
\[
\text{par les } I_{x'_i} \ (i \in I)
\]
Notons que $\mathcal{E}^{\overline{L}}_{L}$ est la limite projective des $\mathcal{E}^{\overline{L}}_{U}$, pour des
modèles réguliers (\uncertain{affines si on y tient}) variables
\[
\text{(24)} \qquad \mathcal{E}^{\overline{L}}_{L} \simeq \varprojlim_{U} \mathcal{E}^{\overline{L}}_{U}
\]
et de même, bien sûr
\[
\text{(25)} \qquad \pi^{\overline{L}}_{L/K} \simeq \varprojlim_{U} \pi^{\overline{L}}_{U/K} .
\]
\marginal{Dans les considérations qui suivent, $K$ (et $X$, $U$) un \ill{} \uncertain{provisoirement} \ill{} \ill{} \ill{} \ill{} [Mais on \uncertain{pourra} \ill{} \ill{} \ill{} \ill{} \ill{} \ill{} \ill{} \uncertain{après coup}] NB \ill{} \ill{} \ill{} \ill{} \ill{} des solutions \uncertain{régulières} \ill{}}
\note{note marginale écrite en oblique dans le coin inférieur gauche, très serrée et en partie barrée ; seuls quelques mots se lisent}
Les \ill{} \ill{} \ill{} les $\mathcal{E}^{\overline{L}}_{U}$ \ill{}
$\mathcal{E}^{\overline{L}}_{L}$ ? Pour \emph{\ill{}} (\emph{\ill{} $V$}) \ill{} du \ill{} de type
\ill{} sur $K$ (correspondant à un « germe de modèle » de $L/K$, \ill{}
\uncertain{voisinage} d'un pt de codim 1 i.e. du pt gén. d'un diviseur), considé\-rons
\struck{\ill{}} le \ill{} fermé $x$ de $S = \operatorname{Spec}(V)$,

\page{117}
et \uncertain{un} $x'$ au dessus de $x$ dans le
normalisé de $S_x$ dans $\overline{L}$, d'où un
groupe de décomposition $N_{x'}$, et un
groupe d'inertie $I_{x'}$, avec
\marginal{\ill{} \uncertain{pourvu} \ill{} d'un « germe » \ill{}}
\[
\text{(26)} \qquad N_{x'} / I_{x'} \simeq \Gamma_{k(x')}
\]
(ext. de type fini de $K$, de degré de \uncertain{transc.} $n-1$, si $n = \operatorname{degtr} L/K$), et
\[
\text{(27)} \qquad I_{x'} \simeq T_{\circ}(\overline{L}) \simeq T_{\infty}(\overline{K})
\]
\marginal{si car. $k(x)$ est \uncertain{nulle} --- en général \ill{} \ill{}…}
\note{l'indice de $T$ est lu $\circ$ dans le premier terme et $\infty$ dans le second, comme sur la page}
(\uncertain{par} Kummer) \struck{Notons ainsi, \ill{}} \add{les \uncertain{constructions}}
\ill{} un \uncertain{système} des \uncertain{choix} \ill{}
\struck{des} modèles $(U, \ldots X \ldots)$. --- Notons que
\[
\text{(28)} \qquad I_{x'} \subset \pi^{\overline{L}}_{L/K} .
\]
En fait, si \uncertain{on} passe aux \uncertain{hensélisés}, \uncertain{posons}
\[
\text{(29)} \quad
\left\{
\begin{aligned}
\tilde{S}_x &= \operatorname{Spec} \tilde{V} , \quad \tilde{V} \text{ hensélisé (pas strict) de } V \\
\tilde{U}_x &= \tilde{S}_x \setminus \{x\} \simeq \operatorname{Spec} \tilde{L}_x \quad (\text{où } \tilde{L}_x \text{ est le corps des fractions de } \tilde{V})
\end{aligned}
\right.
\]
\note{la page porte « $\operatorname{Spec}\tilde{L}_x$ », le second $\tilde L_x$ repassé à l'encre forte}
\struck{\ill{}} donc le choix « effectif » \uncertain{dépend} du
homom. $\mathcal{E}_{\tilde{U}_x} \xrightarrow{\text{ext}} \mathcal{E}_L$,

\page{118}\note{p. 59 de l'auteur}
celui d'un rev. universel \struck{\uncertain{(précis.}} \add{\uncertain{ou} \ill{}} d'une
pt géométrique de $\tilde{U}_x$ --- i.e. de l'une
clôture algébrique de $\tilde{L}_x$ --- \emph{et} d'une
\struck{\ill{}} « domaine de \ill{} » entre $\tilde{U}_x \to U$ \ill{} \uncertain{(Spec L)},
et $\mathcal{E}_L$ défini par $\overline{L}$ : \uncertain{i.e.} \uncertain{compris des situations}
\uncertain{revient} à \ill{}, d'un $\overline{L}$-\ill{} de
\struck{\ill{}} $\tilde{L}_x$ dans $\overline{L}$ --- et ceci \ill{} \ill{}
\marginal{Cf. Hensel \ill{} \ill{} \ill{} \ill{} \ill{} \ill{} \ill{} fait ici \ill{} \ill{} $I_{x'}$ \ill{} \ill{} $\pi^{\overline{L}}_{\tilde{U}_x/k(x)} = \pi^{\overline{L}}_{\tilde{U}_x/\tilde{S}_x}$ !}
donc \uncertain{équivaut} à la donnée d'un
$x'$, \struck{et d'un \ill{}} \uncertain{comme} dessus. \uncertain{Ceci}
posé, on a des \uncertain{isom. canon.}
\[
\text{(30)} \quad
\left\{
\begin{aligned}
N_{x'} &\simeq \mathcal{E}^{\overline{L}}_{\tilde{U}_x} \\
I_{x'} &\simeq \operatorname{Noyau} \text{ de l'homom. canonique} \\
& \qquad \mathcal{E}^{\overline{L}}_{\tilde{U}_x} \longrightarrow \mathcal{E}_{\overline{k(x)}}
\end{aligned}
\right.
\]
\note{au-dessus de la seconde ligne, entre parenthèses et souligné, un mot lu sans certitude « \uncertain{surjectif} »}
($\overline{k(x)}$ : clôture \uncertain{sép.} de $k(x)$ --- \uncertain{induite} par \uncertain{la} clôture alg.
$\overline{L}$ de $\tilde{K}_x$)
\struck{$I_{x'} = \ill{}$}
(\struck{\ill{}}) \uncertain{(dont} (27) est déduit par théorie de Kummer).
\note{l'appel « (27) » est surchargé ; la lecture du chiffre reste incertaine}
\marginal{? (\uncertain{trait vertical double} le long du paragraphe qui suit)}
Supposons que (\uncertain{dans} la situation actuelle,
\uncertain{il existe un} corps de base $K$
\ill{} \uncertain{premier}, \uncertain{\ill{}} $\mathbb{C}$.)
On a \quad $N_{x'} = \operatorname{Norm}_{\mathcal{E}_L}(I_{x'})$, ce qui \uncertain{est lié} \ill{}
à l'idée qu'\uncertain{on devrait avoir} \quad $I_{x'} = \operatorname{Norm}_{\pi_{L/K}}(I_{x'})$

\page{119}
--- à tirer au clair à l'occasion !]

\uncertain{Revenons} aux modèles $U$, fournissant
des quotients $\mathcal{E}^{\overline{L}}_{U}$ de $\mathcal{E}^{\overline{L}}_{L}$, d'où
$N_{x'} \to \mathcal{E}^{\overline{L}}_{U}$, \uncertain{induisant} $I_{x'} \to \mathcal{E}^{\overline{L}}_{U}$,
\struck{Il est clair que} \ill{}
qui \uncertain{n'est autre} que l'homom. \ill{} $\mathcal{E}_{\tilde{U}_x} \to \mathcal{E}_U$
\uncertain{induit} \ill{} par $\tilde{U}_x \to U$. \emph{Si} cet homom.
\uncertain{se} factorise par $\tilde{S}_x$, il est immédiat
\marginal{i.e. si $V$ a un « centre » sur $U$ … \ill{} \ill{} \ill{} \uncertain{local} \ill{} $U$…}
que l'homom. \uncertain{induit} $I_{x'} \to \mathcal{E}^{\overline{L}}_{U}$ est
nul, puisque \struck{\ill{}} la composée de
$\mathcal{E}_{\tilde{U}_x} \to \mathcal{E}_{\tilde{S}_x}$ et $\mathcal{E}_{\tilde{U}_x} \to \mathcal{E}_U$ se factorise par
$\mathcal{E}_{\tilde{U}_x} \to \mathcal{E}_{\tilde{S}_x}$, \ill{} homom. [\uncertain{Avoir} \uncertain{une réciproque} ?
\note{la question est soulignée d'un trait ondulé}
\marginal{Ça devrait \ill{} \ill{} \ill{} \uncertain{éliminant} \ill{} \ill{} \ill{} des « \ill{} » $(0,0)$, $(0,1)$… \ill{} \ill{} \ill{} l'existence \ill{} \ill{} \ill{} \ill{} $K$}
\ill{} ! ]

Quand on dispose d'un
\struck{\ill{}} modèle \emph{propre} \uncertain{dont} $U$ soit un ouvert,
\struck{\ill{}} (par exemple \ill{} \ill{}
\uncertain{en} \ill{} \ill{}) que $V$ \uncertain{domine} \ill{}
\ill{} sur $X$, soit $a$, \struck{et la \ill{}} \add{(\textbf{NB} $a$ \ill{})}
\ill{} de codim.\ 1…) et la condition
\uncertain{nécessaire} \ill{} (\uncertain{c'est-à-dire} que
$\mathcal{E}^{\overline{L}}_{U}$ \uncertain{soit} \uncertain{nul}) est \uncertain{simplement} que $a \in U$.

\page{120}\note{p. 60 de l'auteur}
Ceci posé, $\mathcal{E}^{\overline{L}}_{U}$ se \uncertain{reconstruit} à partir
de $\mathcal{E}^{\overline{L}}_{L}$, comme \uncertain{un} quotient de ce dernier,
en prenant \emph{tous} les $V$ dans $L$
\uncertain{ayant} un centre sur $U$ (i.e.\ \uncertain{tels} \ill{}
$=$ des \uncertain{anneaux} les $V = \mathcal{O}_{U,x}$, où $x$ est
un pt de codim 1 dans $U$), et en
divisant par le \emph{sous-groupe} fermé invariant
engendré par les $I_{x'}$ correspondants.

On peut regarder l'espace \uncertain{loc.} \ill{} ;
(\uncertain{singulier})
\ill{} des modèles \struck{\uncertain{propres plats}} $L$ sur $K$
\[
\text{(31)} \qquad \mathfrak{X}_{L/K} \simeq \varprojlim_{\substack{X \text{ modèle} \\ \text{propre de } L/K}} X
\]
\marginal{\textbf{NB} \ill{} \ill{} \ill{} (31) \ill{} \ill{} \ill{} \ill{} \ill{} \ill{} \ill{} \ill{} \ill{} \ill{} \ill{}…}
\note{marginale oblique très serrée dans le coin inférieur gauche ; seuls les signes « NB » et « (31) » sont sûrs}
\uncertain{(ex. $n \geq 2$)} --- \uncertain{solution}. \struck{\uncertain{Lieu}} \uncertain{Si}
\ill{} on n'a pas un
$x \in \mathfrak{X}$, $k(x)$ est une ext. de $K$ de
degré de transcendance $d \leq n$ --- \ill{} \uncertain{l'appelle}
\uncertain{la} \uncertain{dimension} de $x$ ; $n - d$ est la \emph{codim.} ; il y a un pt et un
\uncertain{seul} de \uncertain{codimension} $0$ --- \uncertain{correspondant} aux pts
gén. de $X$. Les points de codim 1
\note{l'argument se poursuit au-delà de la fin du lot}

\end{document}
